\subsection{不变多项式函数}

设 \(\mathfrak{a}\) 是有限维李代数。按照第 1 节的约定，我们将代数 \(\mathbf{S}(\mathfrak{a}^{*})\)、代数 \(\mathbf{S}(\mathfrak{a})^{*gr}\) 与 \(\mathfrak{a}\) 上多项式函数的代数识别起来。对于所有 \(a \in \mathfrak{a}\)，令 \(\theta(a)\) 为 \(\mathbf{S}(\mathfrak{a})\) 的导子，使 \(\theta(a)x = [a, x]\) 对所有 \(x \in \mathfrak{a}\) 成立。我们知道（第 I 章，§3，第 2 节），\(\theta\) 是 \(\mathfrak{a}\) 在 \(\mathbf{S}(\mathfrak{a})\) 上的表示。令 \(\theta^{*}(\mathfrak{a})\) 为 \(-^{t}\theta(a)\) 在 \(\mathbf{S}(\mathfrak{a}^{*})\) 上的限制。则 \(\theta^{*}\) 是 \(\mathfrak{a}\) 的一个表示。若 \(f \in \mathbf{S}^{n}(\mathfrak{a}^{*})\)，则 \(\theta^{*}(a)f \in \mathbf{S}^{n}(\mathfrak{a}^{*})\)，并且对于 \(x_{1}, \ldots, x_{n} \in \mathfrak{a}\)，

\[
(\theta^ {*} (a) f) (x _ {1}, \dots , x _ {n}) = - \sum_ {1 \leq i \leq n} f (x _ {1}, \dots , x _ {i - 1}, [ a, x _ {i} ], x _ {i + 1}, \dots , x _ {n}).\tag{1}
\]

由 (1) 容易推出 \(\theta^{*}(a)\) 是 \(\mathbf{S}(\mathfrak{a}^{*})\) 的一个导子。属于 \(\mathbf{S}(\mathfrak{a})\)（分别为 \(\mathbf{S}(\mathfrak{a}^{*})\)）且在表示 \(\theta\)（分别为 \(\theta^{*}\)）下不变的元素，称为 \(\mathbf{S}(\mathfrak{a})\)（分别为 \(\mathbf{S}(\mathfrak{a}^{*})\)）的不变元素。

引理 2. 设 \(\rho\) 是 \(\mathfrak{a}\) 的有限维线性表示，\(m\) 是满足 \(\geq 0\) 的整数。函数 \(x \mapsto \operatorname{Tr}(\rho(x)^m)\) 是 \(\mathfrak{a}\) 上的不变多项式函数。

置 \(g(x_{1},\ldots ,x_{m}) = \mathrm{Tr}(\rho (x_{1})\ldots \rho (x_{m}))\)，其中 \(x_{1},\ldots ,x_{m}\in \mathfrak{a}\)。若 \(x\in \mathfrak{a}\)，则有

\[
\begin{array}{l} - (\theta^ {*} (x) g) (x _ {1}, \ldots , x _ {m}) \\ \qquad = \sum_ {1 \leq i \leq m} \mathrm{Tr} (\rho (x _ {1}) \ldots \rho (x _ {i - 1}) [ \rho (x), \rho (x _ {i}) ] \rho (x _ {i + 1}) \ldots \rho (x _ {m})) \\ \qquad = \mathrm{Tr} (\rho (x) \rho (x _ {1}) \ldots \rho (x _ {m})) - \mathrm{Tr} (\rho (x _ {1}) \ldots \rho (x _ {m}) \rho (x)) = 0, \end{array}
\]

所以 \(\theta^{*}(x)g = 0\)。令 \(h\) 为由下式定义的对称多重线性形式：

\[
h (x _ {1}, \dots , x _ {m}) = \frac {1}{m !} \sum_ {\sigma \in \mathfrak {S} _ {m}} g (x _ {\sigma (1)}, \dots , x _ {\sigma (m)}).
\]

对于所有 \(x \in \mathfrak{a}\)，有 \(\theta^{*}(x)h = 0\) 且 \(\operatorname{Tr}(\rho(x)^m) = h(x, \ldots, x)\)，引理得证。

引理 3. 设 E 是有限维 g-模，且 \(x \in E\)。则 x 是 g-模 E 的不变元素，当且仅当对于 g 的每个幂零元 a，都有 \((\exp a_{\mathrm{E}}).x = x\)。

条件显然是必要的。现在假设它满足。设 a 是 g 的幂零元。存在一个整数 n，使得 \(a_{E}^{n}=0\)。对于所有 \(t\in k\)，我们有

\[
0 = \exp (t a _ {\mathrm{E}}). x - x = t a _ {\mathrm{E}} x + \frac {1}{2 !} t ^ {2} a _ {\mathrm{E}} ^ {2} x + \dots + \frac {1}{(n - 1) !} t ^ {n - 1} a _ {\mathrm{E}} ^ {n - 1} x,
\]

所以 \(a_{\mathrm{E}}x = 0\)。但是，李代数 \(\mathfrak{g}\) 由其幂零元生成（§4，第 1 节，命题 1）。因此，\(x\) 是 \(\mathfrak{g}\)-模 E 的不变元素。证毕。

对于任意 \(\xi \in \mathbf{GL}(\mathfrak{g})\)，令 \(\mathbf{S}(\xi)\) 为 \(\mathbf{S}(\mathfrak{g})\) 上延拓 \(\xi\) 的自同构，令 \(\mathbf{S}^*(\xi)\) 表示在 \(\mathbf{S}(\mathfrak{g}^*)\) 上取 \(\mathbf{S}(\xi)\) 的逆变自同构所得的限制。于是 \(\mathbf{S}\) 和 \(\mathbf{S}^*\) 是 \(\mathbf{GL}(\mathfrak{g})\) 的表示。若 \(a\) 是 \(\mathfrak{g}\) 的幂零元，则 \(\theta(a)\) 在 \(\mathbf{S}(\mathfrak{g})\) 上局部幂零，且 \(\mathbf{S}(\exp a) = \exp \theta(a)\)，因此

\[
\mathbf {S} ^ {*} (\exp \operatorname{ad} a) = \exp \theta^ {*} (a).\tag{2}
\]

命题 3. 设 \(f\) 是 \(\mathfrak{g}\) 上的多项式函数。以下条件等价：

\begin{enumerate}
\def\labelenumi{(\roman{enumi})}
\item
  \(f\circ s = f\) 对所有 \(s\in \mathrm{Aut}_e(\mathfrak{g})\)
\item
  \(f\circ s = f\) 对所有 \(s\in \mathrm{Aut}_0(\mathfrak{g})\)
\item
  \(f\) 是不变的。
\end{enumerate}

由公式 (2) 和引理 3 可得 (i) 与 (iii) 等价。由此，通过扩张基域可知 (iii) 蕴含 (ii)。(ii) \(\Longrightarrow\) (i) 显然。

请注意，即使 f 满足命题 3 的条件，f 通常也不具有 \(\operatorname{Aut}(\mathfrak{g})\)-不变性（习题 1 和 2）。

定理 1. 设 \(\mathrm{I}(\mathfrak{g}^*)\) 是 \(\mathfrak{g}\) 上不变多项式函数的代数。设 \(i: \mathbf{S}(\mathfrak{g}^*) \to \mathbf{S}(\mathfrak{h}^*)\) 是限制同态。

\begin{enumerate}
\def\labelenumi{(\roman{enumi})}
\item
  映射 \(i|\mathrm{I}(\mathfrak{g}^*)\) 是从代数 \(\mathrm{I}(\mathfrak{g}^*)\) 到代数 \(\mathbf{S}(\mathfrak{h}^*)^{\mathrm{W}}\) 的同构。
\item
  对任意整数 \(n \geq 0\)，令 \(\mathrm{I}^{n}(\mathfrak{g}^{*})\) 为 \(\mathrm{I}(\mathfrak{g}^{*})\) 中次数为 n 的齐次元素的集合。则 \(\mathrm{I}^{n}(\mathfrak{g}^{*})\) 是 g 上形如 \(x \mapsto \mathrm{Tr}(\rho(x)^{n})\) 的函数的线性组合所成的集合，其中 \(\rho\) 是 g 的有限维线性表示。
\item
  令 \(l = \operatorname{rk}(\mathfrak{g})\)。存在 \(l\) 个 \(\mathrm{I}(\mathfrak{g}^*)\) 中代数独立的齐次元素，它们生成代数 \(\mathrm{I}(\mathfrak{g}^*)\)。
\end{enumerate}

\begin{enumerate}
\def\labelenumi{\alph{enumi})}
\item
  设 \(f \in \mathrm{I}(\mathfrak{g}^{*})\) 且 \(w \in W\)。存在 \(s \in \mathrm{Aut}_{e}(\mathfrak{g}, \mathfrak{h})\)，使得 \(s|_{\mathfrak{h}} = w\)（§2，第 2 节，定理 2 的推论）。由于 f 在 s 下不变（命题 3），\(i(f)\) 在 w 下不变。因此 \(i(\mathrm{I}(\mathfrak{g}^{*})) \subset \mathbf{S}(\mathfrak{h}^{*})^{\mathrm{W}}\)。
\item
  我们证明：若 \(f \in \mathrm{I}(\mathfrak{g}^{*})\) 满足 \(i(f) = 0\)，则 f = 0。必要时扩张基域，可假设 k 是代数闭域。根据命题 3，f 在 \(s(\mathfrak{h})\) 上为零，对所有 \(s \in \operatorname{Aut}_{e}(\mathfrak{g})\) 都成立。因此 f 在 g 的每个嘉当子代数上都为零（第 VII 章，§3，第 2 节，定理 1），特别地，在 g 的正则元集合上为零。但这个集合在 Zariski 拓扑下稠密于 g（第 VII 章，§2，第 2 节）。
\item
  令 n 是满足 \(\geq 0\) 的整数。令 \(L^{n}\) 为 g 上形如 \(x \mapsto \operatorname{Tr}(\rho(x)^{n})\) 的函数的线性组合所成的集合，其中 \(\rho\) 是 g 的有限维线性表示。根据引理 2，\(L^{n} \subset \operatorname{I}^{n}(\mathfrak{g}^{*})\)。于是
\end{enumerate}

\[
i \left(\mathrm{L} ^ {n}\right) \subset i \left(\mathrm{I} ^ {n} \left(\mathfrak {g} ^ {*}\right)\right) \subset \mathbf {S} ^ {n} \left(\mathfrak {h} ^ {*}\right) ^ {\mathrm{W}}.
\]

根据命题 2 的推论 2，\(\mathbf{S}^n (\mathfrak{h}^*)^{\mathrm{W}}\subset i(\mathrm{L}^n)\)。因此 \(i(\mathrm{I}^n (\mathfrak{g}^*)) = \mathbf{S}^n (\mathfrak{h}^*)^{\mathrm{W}}\)，这证明了 (i)；并且 \(i(\mathrm{L}^n) = i(\mathrm{I}^n (\mathfrak{g}^*))\)，所以 \(\mathrm{L}^n = \mathrm{I}^n (\mathfrak{g}^*)\) 根据 \(b\)。这就证明了 (ii)。

\begin{enumerate}
\def\labelenumi{\alph{enumi})}
\setcounter{enumi}{3}
\tightlist
\item
  断言 (iii) 由 (i) 以及第 V 章，§5，第 3 节，定理 3 得出。
\end{enumerate}

推论 1. 假设 \(\mathfrak{g}\) 是单的。令 \(m_1, \ldots, m_l\) 为 \(\mathfrak{g}\) 的 Weyl 群的指数。存在 \(\mathrm{P}_1, \ldots, \mathrm{P}_l\)，它们是 \(\mathrm{I}(\mathfrak{g}^*)\) 中的齐次元素，次数分别为

\[
m _ {1} + 1, \dots , m _ {l} + 1,
\]

它们代数无关，并生成代数 \(\mathrm{I}(\mathfrak{g}^*)\)。

这由定理 2 (i) 以及第五章，§6，第 2 项，命题 3 推出。

推论 2. 令 B 为 R 的一个基，\(R_{+}\)（分别为 \(R_{-}\)）为相对于 B 的 \((\mathfrak{g},\mathfrak{h})\) 的正（分别为负）根，\(n_{+}=\sum_{\alpha\in R_{+}}g^{\alpha}\)，\(n_{-}=\sum_{\alpha\in R_{-}}g^{\alpha}\)，\(\mathbf{S}(\mathfrak{h})\) 为 h 的对称代数，而 J 为 \(\mathbf{S}(\mathfrak{g})\) 中由 \(n_{+}\cup n_{-}\) 生成的理想。

\begin{enumerate}
\def\labelenumi{(\roman{enumi})}
\item
  \(\mathbf{S}(\mathfrak{g}) = \mathbf{S}(\mathfrak{h}) \oplus \mathbf{J}.\)
\item
令 j 为从代数 \(\mathbf{S}(\mathfrak{g})\) 到代数 \(\mathbf{S}(\mathfrak{h})\) 的同态，它由前述 \(\mathbf{S}(\mathfrak{g})\) 的分解定义。令 \(\mathrm{I}(\mathfrak{g})\) 为 \(\mathbf{S}(\mathfrak{g})\) 的不变元素所成的集合。令 \(\mathbf{S}(\mathfrak{h})^{\mathrm{W}}\) 为 \(\mathbf{S}(\mathfrak{h})\) 中在 W 的作用下不变的元素所成的集合。则 \(j|I(\mathfrak{g})\) 是从 \(\mathrm{I}(\mathfrak{g})\) 到 \(\mathbf{S}(\mathfrak{h})^{\mathrm{W}}\) 的同构。
\end{enumerate}

断言 (i) 显然成立。Killing 形式定义了从向量空间 \(\mathfrak{g}^*\) 到向量空间 \(\mathfrak{g}\) 的同构，并将其扩张为同构 \(\xi\)，从 \(\mathfrak{g}\)-模 \(\mathbf{S}(\mathfrak{g}^*)\) 到 \(\mathfrak{g}\)-模 \(\mathbf{S}(\mathfrak{g})\)。我们有 \(\xi (\mathrm{I}(\mathfrak{g}^*)) = \mathrm{I}(\mathfrak{g})\)。相对于 Killing 形式，\(\mathfrak{h}\) 的正交补是 \(\mathfrak{n}_{+} + \mathfrak{n}_{-}\)（§2，第 2 项，命题 1）。若将 \(\mathfrak{h}^*\) 认同为 \(\mathfrak{n}_{+} + \mathfrak{n}_{-}\) 在 \(\mathfrak{g}^*\) 中的正交补，则 \(\xi (\mathfrak{h}^*) = \mathfrak{h}\)，从而 \(\xi (\mathbf{S}(\mathfrak{h}^*)) = \mathbf{S}(\mathfrak{h})\)，且 \(\xi (\mathbf{S}(\mathfrak{h}^*)^{\mathrm{W}}) = \mathbf{S}(\mathfrak{h})^{\mathrm{W}}\)。最后，\(\xi^{-1}(\mathrm{J})\) 是 \(\mathfrak{g}\) 上在 \(\mathfrak{h}\) 上消失的多项式函数所成的集合。这证明 \(\xi\) 将定理 1 中的同态 \(i\) 变为推论 2 中的同态 \(j\)。因此，断言 (ii) 由定理 1 (i) 得出。

命题 4. 设 \(\mathfrak{a}\) 为半单李代数，\(l\) 为其秩。令 I（分别为 I'）为 \(\mathbf{S}(\mathfrak{a}^{*})\)（分别为 \(\mathbf{S}(\mathfrak{a})\)）中在由 \(\mathfrak{a}\) 的伴随表示诱导的表示下不变的元素所成的集合。令 Z 为 \(\mathfrak{a}\) 的包络代数的中心。

\begin{enumerate}
\def\labelenumi{(\roman{enumi})}
\item
  I 和 \(\mathrm{I}'\) 是分次多项式代数（第五章，§5，第 1 项），其超越次数为 \(l\)。
\item
  \(Z\) 同构于 \(l\) 个以 \(k\) 为系数的不定元的多项式代数。
\end{enumerate}

\(\mathfrak{a}\) 的包络代数的典则滤过在 \(Z\) 上诱导出一个滤过。根据第一章，§2，第 7 项，定理 1 及 p.~25，gr Z 同构于 I'。鉴于《交换代数》第三章，§2，第 9 项，命题 10，便可推出 (i) \(\Longrightarrow\) (ii)。

另一方面，定理 1 及其推论 2 表明，当 \(\mathfrak{a}\) 可分裂时，(i) 成立。一般情形根据下述引理归约到该情形：

引理 4。\(^{2}\) 令 \(A = \bigoplus_{n \geq 0} A^{n}\) 为一个分次 k-代数，\(k'\) 为 k 的一个扩域，且 \(A' = A \otimes_{k} k'\)。假设 \(A'\) 是 \(k'\) 上的分次多项式代数。则 A 是 k 上的分次多项式代数。

我们有 \(A^{\prime0}=k'\)，所以 \(A^{0}=k\)。置 \(A_{+}=\bigoplus_{n\geq1}A^{n}\) 且 \(P=A_{+}/A_{+}^{2}\)。于是 P 是一个分次向量空间，并且存在一个次数为零的分次线性映射 \(f:P\to A_{+}\)，使得它与典则投影 \(A_{+} \rightarrow P\) 的复合是 P 上的恒等映射。给 \(\mathbf{S}(P)\) 赋予由 P 的分次结构诱导的分次结构（《代数》，第三章，p.~506）。延拓 f 的 k-代数同态 \(g : \mathbf{S}(P) \rightarrow A\)（《代数》，第三章，p.~497）是一个次数为 0 的分次同态；立即对次数作归纳可见，g 是满射。

引理 5。A 是分次多项式代数，当且仅当 P 是有限维的且 g 是双射。

若 P 是有限维的，则 \(\mathbf{S}(\mathrm{P})\) 显然是分次多项式代数；若 g 为双射，A 也是如此。反过来，假设 A 由代数无关的齐次元素 \(x_{1},\ldots,x_{m}\) 生成，其次数为 \(d_{1},\ldots,d_{m}\)。令 \(\bar{x}_{i}\) 为 \(x_{i}\) 在 P 中的像。显然 \(\bar{x}_{i}\) 构成 P 的一个基；由于 \(\bar{x}_{i}\) 的次数是 \(d_{i}\)，所以 \(\mathbf{S}(\mathrm{P})\) 与 A 同构；特别地，对所有 n 都有 \(\dim\mathbf{S}(\mathrm{P})^{n}=\dim\mathrm{A}^{n}\)。由于 g 是满射，它必为双射。

引理 4 现在立即可得。事实上，将引理 5 应用于 \(k'\)-代数 \(\mathrm{A}'\)，可知 \(g \otimes 1: \mathbf{S}(\mathrm{P}) \otimes k' \to \mathrm{A} \otimes k'\) 是双射，因而 \(g\) 也是双射。

命题 5。沿用命题 4 的记号，并记 p 为 \(\mathbf{S}(\mathfrak{a}^{*})\) 中由 I 的次数 \(\geq 1\) 的齐次元素生成的理想。令 \(x \in a\)。则 x 是幂零的，当且仅当 \(f(x) = 0\) 对所有 \(x \in p\) 都成立。³

如有必要可扩张基域，设 \(\mathfrak{a} = \mathfrak{g}\) 为可分裂的。假设 \(x\) 是幂零的。对于任意有限维线性表示 \(\rho\)（属于 \(\mathfrak{g}\)）及任意整数 \(n \geq 1\)，有 \(\operatorname{Tr}(\rho(x)^n) = 0\)，所以有 \(f(x) = 0\)，其中 \(f \in \mathrm{I}(\mathfrak{g}^*)\) 为次数 \(\geq 1\) 的齐次元素（定理 1 (ii)）；此外有 \(f(x) = 0\) 对所有 \(f \in \mathfrak{p}\) 成立。反之，若有 \(f(x) = 0\) 对所有 \(f \in \mathfrak{p}\) 成立，则有 \(\operatorname{Tr}((\operatorname{ad} x)^n) = 0\) 对所有 \(n \geq 1\) 成立（定理 1 (ii)），所以 \(x\) 是幂零的。

*备注 1) 令 \(\mathrm{P}_1, \ldots, \mathrm{P}_l\) 为 I 中生成代数 I 的代数无关齐次元素。则 \((\mathrm{P}_1, \ldots, \mathrm{P}_l)\) 是 \(\mathbf{S}(\mathfrak{a}^*)\) -正则序列（第五章，§5，第 5 项）。事实上，如有必要可扩张基域，设 \(\mathfrak{a} = \mathfrak{g}\) 为可分裂的。现在令 \(\mathrm{N} = \dim \mathfrak{g}\)，并令

\[
\left(\mathrm{Q} _ {1}, \dots , \mathrm{Q} _ {\mathrm{N} - l}\right)
\]

为 \(\mathfrak{h}\) 在 \(\mathfrak{g}^*\) 中的正交补的一个基。令 \(\mathfrak{m}\) 为 \(\mathbf{S}(\mathfrak{g}^*)\) 中由 \(\mathrm{P}_1, \ldots, \mathrm{P}_l, \mathrm{Q}_1, \ldots, \mathrm{Q}_{\mathrm{N}-l}\) 生成的理想。则 \(\mathbf{S}(\mathfrak{g}^*)\mathfrak{m}\) 同构于 \(\mathbf{S}(\mathfrak{h}^*)/\mathrm{J}\)，其中 \(\mathrm{J}\) 是 \(\mathbf{S}(\mathfrak{h}^*)\) 中由 \(i(\mathrm{P}_1), \ldots, i(\mathrm{P}_l)\) 生成的理想。根据定理 1 以及第五章，§5，第 2 项，定理 2，\(\mathbf{S}(\mathfrak{h}^*)/\mathrm{J}\) 是有限维向量空间，因而 \(\mathbf{S}(\mathfrak{g}^*)/\mathfrak{m}\) 也是有限维的。根据《交换代数》中的一个结果，便可推出 \((\mathrm{P}_1, \ldots, \mathrm{P}_l, \mathrm{Q}_1, \ldots, \mathrm{Q}_{\mathrm{N}-l})\) 是 \(\mathbf{S}(\mathfrak{g}^*)\) -正则序列，从而 \((\mathrm{P}_1, \ldots, \mathrm{P}_l)\) 也如此。

\begin{enumerate}
\def\labelenumi{\arabic{enumi})}
\setcounter{enumi}{1}
\tightlist
\item
  代数 \(\mathbf{S}(\mathfrak{a}^*)\) 是 I 上的分次自由模。事实上，这由命题 4、备注 1 以及第五章，§5，第 5 项，引理 5 得出。*
\end{enumerate}

